% == == == == == == == == == == == == == == == == == == == == == == == == == ==
    == == == == == == == == == == == ==
    = % Section 7 : Network Logic Diagram(NLD) % == == == == == == == == == ==
      == == == == == == == == == == == == == == == == == == == == == == == == ==
      == == ==
    =

\section {
  Network Logic Diagram
}
\label {
sec:
  nld_sec
}
\sectionauthors{CM, JH} % Reviewed by Jeremie,
    do we think we can add the full nld just broken up here instead of section
        by section,
    so clint can see the flow
\subsection{NLD Overview}

            The Network Logic Diagram illustrates the chronological sequence and
                logical dependencies of all project activities
                    .It links every task from the WBS to its
                        necessary predecessors and
                            successors to visualize the workflow
                    .This structure allows the team to identify the Critical
                        Path which is the longest chain of dependent tasks that
                            determines the minimum project duration.Isolating
                                these critical tasks helps the team manage
                                    schedule risks and monitor float
                                        within non -
            critical activities.The complete diagram includes 84 activities and
        89 dependencies organized by phase.

\subsection{NLD Structure}

            The diagrams use the Activity -
            on -
            Node(AON) format.Each activity node displays information as follows
    :

\begin{figure}[h !]
    \centering
    \texttt {
  ES | EF
}
\\[3pt]
    \fbox {
  \parbox{2.5cm} {\centering Task Name \\ (Duration) }
}
\caption { Activity - on - Node(AON) Structure }
\label {
fig:
  nld_structure
}
\end { figure }

\noindent\textbf { Node Components: }
\begin { itemize }
    \item \textbf{ES (Early Start)}: Earliest day the task can begin, based on predecessor completion
    \item \textbf{EF (Early Finish)}: Earliest day the task can complete (ES + Duration)
    \item \textbf{Task Name}: Descriptive activity name
    \item \textbf{Duration}: Number of working days for the task
\end{
      itemize
    }

    \noindent\textbf { Arrow Styles: }
    \begin { itemize }
    \item \textbf{Red arrows} : Critical path dependencies;
    any delay in these connections will extend the project duration
    \item \textbf{Gray arrows}
        : Non -
          critical dependencies with available float / slack time
\end {
      itemize
    }

\noindent\textbf{Phase Color Coding:}
Nodes are colored by project phase for visual clarity: Research (blue), Design (yellow), Build (orange), Test (green), and Close (pink).

All arrows indicate Finish-to-Start (FS) dependencies, where the successor cannot start until the predecessor completes.

\subsection{NLD Diagram}

The network diagrams below show task dependencies organized by phase, with critical path indicated by red arrows.

% NLD Legend explaining AON format (auto-generated)
% NLD Legend - AON Node Format Explanation % Generated : 2025 - 12 -
        06 16 : 42

\begin{figure}[H]
\centering
\begin{tikzpicture}[node distance = 1cm and 2cm,
                     legend_node /.style =
                         {rectangle, draw, rounded corners = 3pt,
                          minimum width = 2.5cm, minimum height = 0.9cm,
                          align = center, font =\scriptsize},
                     timing /.style = {font =\tiny, inner sep = 1pt},
                     label_node /.style = {font =\scriptsize, text width = 4cm,
                                           align = left},
                     arrow_label /.style = {font =\scriptsize},
                     critical_arrow /.style = {-{Stealth[length = 6pt]}, red,
                                               line width = 1.5pt},
                     regular_arrow /.style = {-{Stealth[length = 6pt]},
                                              gray !70, thick},
]

            % Example AON node with ES |
    EF
  \node[legend_node, fill = researchblue !30](regular){Task Name\\(5d)};
\node[timing, above = 3pt of regular]{0 | 5};
\node[label_node, right = 1.5cm of regular](reg_label){
    \textbf{AON Node Format}\\[2pt]
    \textbf{Above : } ES | EF\\
    \quad ES = Early Start\\
    \quad EF = Early Finish\\[2pt]
    \textbf{Inside : } Task name, duration};

% Phase colors
  \node[legend_node, fill = designyellow !30, below = 1.5cm of regular](
      phase_ex){Design Task\\(3d)};
\node[timing, above = 3pt of phase_ex]{5 | 8};
\node[label_node, right = 1.5cm of phase_ex](phase_label){
    \textbf{Phase Colors}\\[2pt]
    \fcolorbox{black} {
    researchblue !30} {\phantom{\rule{0.8em} {0.8em}}} Research\\
    \fcolorbox{black} {
    designyellow !30} {\phantom{\rule{0.8em} {0.8em}}} Design\\
    \fcolorbox{black} {buildorange !30} {\phantom{\rule{0.8em} {0.8em}}} Build\\
    \fcolorbox{black} {testgreen !30} {\phantom{\rule{0.8em} {0.8em}}} Test\\
    \fcolorbox{black} {closepink !30} {\phantom{\rule{0.8em} {0.8em}}} Close};

% Arrow styles
  \node[below = 2.5cm of phase_ex, xshift = -0.5cm](arrow_start1){};
\node[right = 2.5cm of arrow_start1](arrow_end1){};
\draw[critical_arrow](arrow_start1)--(arrow_end1);
\node[arrow_label, right = 0.3cm of arrow_end1]{Critical Path Dependency};

\node[below = 0.6cm of arrow_start1](arrow_start2){};
\node[right = 2.5cm of arrow_start2](arrow_end2){};
\draw[regular_arrow](arrow_start2)--(arrow_end2);
\node[arrow_label, right = 0.3cm of arrow_end2]{Regular Dependency};

\end { tikzpicture }
\caption { NLD Node Format Legend(Activity - on - Node) }
\label {
fig:
  nld_legend
}
\end { figure }

\noindent\textbf { Reading the Diagrams: }
\begin{itemize}[nosep]
    \item Each node represents a project activity / task
    \item \textbf{ES}(Early Start)
    : Earliest day the task can begin
    \item \textbf{EF}(Early Finish)
    : Earliest day the task can complete
    \item \textbf{Red arrows} : Critical path,
    any delay extends project duration
    \item \textbf{Gray arrows} : Non - critical dependencies with float / slack
                                            time
\end {
  itemize
}


% Native LaTeX NLD using TikZ (auto-generated)
% NLD Overview - Phase Connections % Generated : 2025 - 12 -
    06 16 : 42

\subsection{Network Logic Diagram Overview}

    The NLD shows task dependencies and the critical path through the
        project.Tasks are organized by phase,
    with arrows indicating predecessor - successor relationships.

\vspace {
  0.5em
}

\begin{figure}[H]
\centering
\resizebox{\textwidth} { ! }
{
  %
\begin{tikzpicture}[node distance = 2cm,
                     phase /.style = {rectangle, draw, rounded corners = 5pt,
                                      minimum width = 3cm,
                                      minimum height = 1.2cm, align = center,
                                      font =\small\bfseries},
                     arrow /.style = {-{Stealth[length = 8pt]}, thick,
                                      gray !70},
  ]

  \node[phase, fill = researchblue](research){Research\\Phase};
  \node[phase, fill = designyellow, right = of research](design){Design\\Phase};
  \node[phase, fill = buildorange, right = of design](build){Build\\Phase};
  \node[phase, fill = testgreen, right = of build](test){Test\\Phase};
  \node[phase, fill = closepink, right = of test](close){Close\\Phase};

  % Phase connections
  \draw[arrow](research)--(design);
  \draw[arrow](design)--(build);
  \draw[arrow](build)--(test);
  \draw[arrow](test)--(close);

  \end { tikzpicture }
}
%
\caption { Project Phase Flow }
\label {
fig:
  nld_overview
}
\end { figure }

\begin{table}[H]
\centering
\caption { Phase Statistics }
\label {
tab:
  nld_stats
}
\begin{tabular} {
  @{} lrrr @{}
}
\toprule
\textbf{Phase} & \textbf{Tasks} & \textbf{Duration} & \textbf {
  Critical
}
\\
\midrule Research & 13 & 56d & 9 \ Design & 24 & 103d & 9 \ Build & 20 & 94d &
    16 \ Test & 15 & 62d & 11 \ Close & 7 & 24d & 3 \\
\midrule
\textbf{Total} & \textbf{79} & \textbf{339d} & \textbf {
  48
}
\\
\bottomrule
\end { tabular }
\end { table }


% Critical Path Analysis - multi-row by phase (auto-generated)
\input{
  nld / nld_critical_path
}

\subsection{Phase - Specific Network Diagrams}

The following diagrams show detailed task dependencies within each project
    phase.Red arrows indicate critical path dependencies;
gray arrows indicate non - critical dependencies with available float time.ES |
    EF values are displayed above each node.

        % Research Phase NLD(auto - generated)
\input{nld / nld_research}

        % Design Phase NLD(auto - generated)
\input{nld / nld_design}

        % Build Phase NLD(auto - generated)
\input{nld / nld_build}

        % Test Phase NLD(auto - generated)
\input{nld / nld_test}

        % Close Phase NLD(auto - generated)
\input {
  nld / nld_close
}

\subsection{Critical Path Method (CPM)}

The critical path was calculated using the Critical Path Method, which involves a forward pass to calculate Early Start (ES) and Early Finish (EF) times, followed by a backward pass to calculate Late Start (LS) and Late Finish (LF) times. Total Float for each task is calculated as LS - ES.

Tasks with zero float form the critical path and any delay in these tasks will delay the project completion date. The complete CPM calculation table showing ES, EF, LS, LF, and Float for all 84 activities is provided in \hyperref[app:cpm_table]{Appendix C.1}.

\subsection{NLD Validation}

The NLD has been validated against the WBS to ensure:

\begin{
  itemize
}
\item All 84 WBS leaf activities appear in the NLD
    \item Every activity(except START / END) has at least one input and output
    \item No circular dependencies exist
    \item The critical path is clearly identifiable
    \item Activity durations match WBS estimates
\end {
  itemize
}

\subsection{Dependency Chains}

The following dependency chains represent critical project flows :

\begin {
  enumerate
}
\item \textbf{Hardware Development Chain}
    : Requirements $\rightarrow$ Component Selection $\rightarrow$ Mechanical
          Design $\rightarrow$ Chassis Fabrication $\rightarrow$ System
              Integration

    \item \textbf{Software Development Chain}
    : Requirements $\rightarrow$ Software Architecture $\rightarrow$ SLAM
          Development $\rightarrow$ Navigation Development $\rightarrow$ System
              Integration

    \item \textbf{Validation Chain}
    : System Integration $\rightarrow$ Unit Testing $\rightarrow$ Integration
          Testing $\rightarrow$ System Validation $\rightarrow$ Acceptance
              Testing
\end{enumerate}

      These chains inform resource allocation and risk management priorities.
